\textbf{COLUNI 2018 -- Questão 07}\par\medskip Leia os fragmentos dos contos:

\par

“Mas de que cor eram os olhos de minha mãe? Eu sabia, desde aquela época, que a mãe inventava esse e outros jogos para distrair a nossa fome. E a nossa fome se distraía.” (Olhos d'água, p. 17)

\par

“Estava, porém, chegando à conclusão de que trabalho como o dela não resolvia nada. Mas o que fazer? Se parasse, a fome viria mais rápida e voraz ainda.” (Zaíta esqueceu de guardar os brinquedos, p. 75)

\par

A fome, devido às dificuldades de sobrevivência enfrentadas pelos personagens, é um assunto que vem à tona nos contos de Conceição Evaristo e é, frequentemente, referida como algo que possui características animadas. No livro de Carolina Maria de Jesus, a fome também é recorrente.

\par

Assinale a alternativa em que Carolina Maria de Jesus usa a mesma estratégia linguística que a usada, nos fragmentos acima, para se referir à fome:\begin{enumerate}
\item “Os meus filhos estão sempre com fome. Quando eles passam muita fome eles não são exigentes no
      paladar.” (p. 30)
\item “Só que o Cesar da atualidade supera o Cesar do passado. Os outros era perseguido pela fé. E nós, pela
      fome!” (p. 146)
\item “Eu suicidando-me é por deficiencia de alimentação no estomago. E por infelicidade eu amanheci com
      fome.” (p. 99)
\item “Acendeu o fogo e assou a carne. A fome era tanta que ele não poude deixar assar a carne.” (p. 40)\end{enumerate}
